\section{A quantitative guide to application regimes}\label{sec:scenarios}

\Cref{sec:applications} lists application areas. Here the dimensionless groups that decide the heat law are evaluated for six representative scenarios, so that the reader can see at a glance which of the models of this review is relevant. The table is generated by \path{scripts/table_scenarios.py} from the rounded handbook inputs of \cref{tab:materials} and a lag $\tau_q=1$\,ps; $\delta_T=\sqrt{2\chi/\omega}$ is the thermal penetration depth and $\lambda=c_1/f$ the acoustic wavelength.

\begin{table}[ht]\centering\scriptsize
\caption{Dimensionless groups for six scenarios ($T_0=300$\,K, $\tau_q=1$\,ps). The assessment column applies the rule: Fourier is adequate if $\omega\tau_q<0.1$ and $\omega/\omega_*<0.1$; a lag law may matter if $\omega\tau_q<1.5$ but then the kinetic validity of the model must be checked; otherwise the continuum heat law itself is questionable.}\label{tab:scenarios}
\setlength{\tabcolsep}{3pt}\begin{tabular}{@{}p{3.0cm}lrrrrrp{2.3cm}@{}}\toprule
Scenario & Material & $f$ (Hz) & $\omega/\omega_*$ & $\omega\tau_q$ & $\delta_T$ (m) & $\lambda$ (m) & Assessment\\\midrule
Ultrasonic testing & Aluminium & $1.0\times10^{7}$ & $1.6\times10^{-4}$ & $6.3\times10^{-5}$ & $1.8\times10^{-6}$ & $6.2\times10^{-4}$ & Fourier adequate\\
GHz acoustics & Aluminium & $1.0\times10^{9}$ & $1.6\times10^{-2}$ & $6.3\times10^{-3}$ & $1.8\times10^{-7}$ & $6.2\times10^{-6}$ & Fourier adequate\\
Picosecond ultrasonics & Aluminium & $1.0\times10^{11}$ & $1.6\times10^{0}$ & $6.3\times10^{-1}$ & $1.8\times10^{-8}$ & $6.2\times10^{-8}$ & lag may matter; check kinetics\\
Silicon beam, $h=1\,\mu$m, at $f_Z$ & Silicon & $1.4\times10^{8}$ & $1.1\times10^{-3}$ & $9.1\times10^{-4}$ & $4.5\times10^{-7}$ & $6.1\times10^{-5}$ & Fourier adequate\\
Silicon film, $h=20$\,nm, at $f_Z$ & Silicon & $3.6\times10^{11}$ & $2.7\times10^{0}$ & $2.3\times10^{0}$ & $9.0\times10^{-9}$ & $2.5\times10^{-8}$ & continuum law questionable\\
Copper film, $h=50$\,nm, at $f_Z$ & Copper & $7.3\times10^{10}$ & $2.6\times10^{0}$ & $4.6\times10^{-1}$ & $2.3\times10^{-8}$ & $6.2\times10^{-8}$ & lag may matter; check kinetics\\
\bottomrule\end{tabular}
\end{table}

\subsection{Reading the table}
\begin{itemize}
\item \emph{Ultrasonic testing and GHz acoustics (aluminium).} With $\omega\tau_q$ between $6\times10^{-5}$ and $6\times10^{-3}$ and $\omega/\omega_*$ between $1.6\times10^{-4}$ and $1.6\times10^{-2}$, classical coupled thermoelasticity is adequate. The thermal penetration depth ($1.8\,\mu$m at $10$\,MHz, $0.18\,\mu$m at $1$\,GHz) is much smaller than the acoustic wavelength ($620\,\mu$m and $6.2\,\mu$m), so the wave sees the temperature field as a thin boundary layer at interfaces and as a small bulk damping.
\item \emph{Picosecond ultrasonics (aluminium, 100\,GHz).} Here $\omega/\omega_*=1.6$ and $\omega\tau_q=0.63$ and the wavelength ($62$\,nm) is comparable with the thermal penetration depth ($18$\,nm). This is the regime of the regime map (\cref{sec:synthesis}) where lag and coupling effects could be visible, and also where carrier mean free paths are comparable with the wavelength, so the interpretation of any deviation from the classical result requires a kinetic or atomistic cross-check.
\item \emph{Silicon beam, $h=1\,\mu$m.} At the Zener frequency $f_Z=1.4\times10^{8}$\,Hz, $\omega\tau_q\approx10^{-3}$ and the Fourier law with the exact flexural result \cref{eq:LR} is adequate. The thermal penetration depth ($0.45\,\mu$m) is a sizeable fraction of the thickness, which is why the damping peaks here.
\item \emph{Silicon film $20$\,nm and copper film $50$\,nm.} At the thermal relaxation frequency of the film, $\omega\tau_q$ is $2.3$ and $0.46$ and $\omega/\omega_*$ is $2.7$ and $2.6$. These are the regimes in which a continuum heat law with bulk parameters has to be treated with caution: the film thickness is comparable to or below the carrier mean free path, boundary scattering reduces the effective conductivity, and the relaxation frequency of the continuum model is not reliable.
\end{itemize}

\subsection{Consequences for the choice of theory}
\begin{enumerate}
\item The classical coupled theory suffices for essentially all ultrasonic, MEMS-scale-above-micron and geophysical-scale applications, and it is the one to use as a baseline in any case.
\item Lag models can only matter in the lower right-hand corner of \cref{tab:scenarios}: frequencies above about $10^{10}$\,Hz or thicknesses below about $100$\,nm. In that corner the continuum assumptions are also stressed, so any claim of a lag effect needs an independent test.
\item The Guyer--Krumhansl and nonlocal models (\cref{sec:kinetic,sec:size}) are the first refinement once size effects are suspected, because they carry a length scale; lag models do not.
\item The electromechanical and magnetic couplings (\cref{sec:multi}) change the stiffness at all frequencies, and are therefore relevant over the whole table, even where the heat law is not.
\end{enumerate}
